\documentclass[10pt]{article}

Така получаваме, че
\[
\nabla \left( \frac{1}{\|x\|} \right) = \left( -\frac{x_1}{r^3},\; -\frac{x_2}{r^3},\; -\frac{x_3}{r^3} \right) = F(x).
\]

Следователно функцията
\[
u(x) = \frac{1}{\|x\|}
\]
е потенциал на полето $F(x) = -\dfrac{x}{\|x\|^3}$.

Тъй като потенциалът е определен с точност до константа, общият вид на потенциала е
\[
u(x) = \frac{1}{\|x\|} + C,
\]
където $C \in \mathbb{R}$ е произволна константа. Обикновено се избира $C = 0$, така че $\displaystyle \lim_{\|x\|\to\infty} u(x) = 0$.

\[
\boxed{u(x) = \dfrac{1}{\|x\|}}
\]

\textbf{Задача 12(2,3) }Функцията $f : [-\pi/2, \pi/2] \to \mathbb{R}$, зададена по-долу, интегруема ли е по Риман? Обосновете отговора си.
\[
f(x) = 
\begin{cases}
1, & \text{ako } \sin(1/x) > 0,\; x \ne 0; \\
-1, & \text{ako } \sin(1/x) < 0,\; x \ne 0; \\
0, & \text{ako } \sin(1/x) = 0,\; x \ne 0; \\
0, & \text{ako } x = 0.
\end{cases}
\]

\textbf{Решение.} 
Функцията $f$ е дефинирана в затворения интервал $[-\pi/2,\pi/2]$ и е ограничена, тъй като $|f(x)| \le 1$ за всяко $x$.

Ще използваме критерия на Лебег за интегруемост по Риман:

\vspace{0.2cm}
\noindent
\textit{Една ограничена функция е интегруема по Риман в даден интервал точно тогава, когато множеството от точките ѝ на прекъсване има Лебегова мярка нула.}

\vspace{0.2cm}
\noindent
Функцията $\sin(1/x)$ е непрекъсната за $x \ne 0$. Следователно $f$ е непрекъсната във всяка точка $x \ne 0$, в която $\sin(1/x) \ne 0$. Прекъсване може да има само:
\begin{itemize}
    \item в точката $x = 0$;
    \item в точките, където $\sin(1/x) = 0$, т.е. $1/x = k\pi$, $k \in \mathbb{Z}$.
\end{itemize}

От $\sin(1/x) = 0$ получаваме $x = \dfrac{1}{k\pi}$, $k \in \mathbb{Z}$. В интервала $[-\pi/2,\pi/2]$ това са точките
\[
x_k = \frac{1}{k\pi}, \quad k \in \mathbb{Z},\; |k| \ge 1,
\]
както и $x = 0$, което съответства на граничния случай $k \to \infty$.

Множеството от точките на прекъсване на $f$ е
\[
D = \{0\} \cup \left\{ \frac{1}{k\pi} \;:\; k \in \mathbb{Z},\; |k| \ge 1 \right\}.
\]

Това множество е \textit{изброимо} – обединение на една точка и изброимо много точки. Всяка точка има Лебегова мярка нула, а изброимото обединение на множества с мярка нула също има мярка нула.

Следователно $D$ е \textit{пренебрежимо по Лебег} (има мярка нула).

Тъй като $f$ е ограничена и множеството от точките ѝ на прекъсване е с мярка нула, от \textit{критерия на Лебег} следва, че $f$ е \textit{интегруема по Риман} в $[-\pi/2,\pi/2]$.

\[
\boxed{\text{Функцията е интегруема по Риман.}}
\]



\section*{3}

\textbf{Задача 1(1) }Докажете, че контурът на обединението на две множества се съдържа в обединението на техните контури: $\partial(A \cup B) \subset \partial A \cup \partial B$. 

\textbf{Задача 2(4,5) }Каква е стойността на $\operatorname{rot}(\operatorname{grad} f)$ и защо?

\textbf{Задача 3(5,6) }Разгледайте функцията $f(x) = e^{3\|x\|}\langle x,a\rangle$, където $x = (x_1,x_2,x_3) \in \mathbb{R}^3$ и $a$ е векторът $(2,1,3)$. Пресметнете $\operatorname{grad} f$ и $\operatorname{div}(\operatorname{grad} f)$.

\section*{2}

\textbf{Задача 1(-) }Да разгледаме гладкото векторно поле $F(x) = x \sin\|x\|$, където $x = (x_1,x_2,x_3) \in \mathbb{R}^3$. Пресметнете $\operatorname{div} F(x)$ и $\operatorname{rot} F(x)$. 

\textbf{Задача 2(7-8) }Да разгледаме гладкото векторно поле $F(x) = x e^{\|x\|}$, където $x = (x_1,x_2,x_3) \in \mathbb{R}^3$. Пресметнете $\operatorname{div} F(x)$ и $\operatorname{rot} F(x)$.

\textbf{Задача 3(6) }(б) Пресметнете $\operatorname{div} \operatorname{rot} F$. 

\textbf{Задача 4(5,6) }Формулирайте и докажете формулата на Лайбниц за диференциране по параметър под знака на интеграла. 

\textbf{Задача 5(-) }Представете тройния интеграл $\iiint_K g(x,y,z) \,dxdydz$ като повторен. Тук $g$ е непрекъсната функция, дефинирана върху тялото на Вивиани $K = \{(x,y,z)\in\mathbb{R}^3: x^2+y^2+z^2 \le 1,\; x^2+y^2 \le x\}$. 

\textbf{Задача 6(7) }Нека $\varphi_1$ и $\varphi_2$ ($\varphi_1<\varphi_2$) са дължините на два меридиана, а $\psi_1$ и $\psi_2$ ($\psi_1<\psi_2$) са ширините на два паралела (в радиани) върху сфера с радиус $R$. Пресметнете площта на частта от сферата, заключена между тези два паралела и два меридиана. 
\end{document}
